\documentclass[12pt, oneside, a4paper, brazil]{abntex2}
% \usepackage[utf8]{inputenc}
\usepackage[brazil]{babel}
% \usepackage[T1]{fontenc}
\usepackage{fontspec}
\usepackage{setspace}
\usepackage{graphicx}
\usepackage{scalefnt}
\usepackage{float}
\usepackage[a4paper, left=3cm, right=2.5cm, top=3cm, bottom = 2.5cm]{geometry}
\usepackage{hyperref}
\usepackage{listings}
\usepackage{color}
\usepackage{titlesec}
%\usepackage[alf]{abntex2cite}
% \usepackage{times}
\usepackage[scaled]{helvet}
\usepackage{fancyhdr}
\usepackage{lmodern}
\usepackage{indentfirst}

%%% Definição de variáveis. %%%
% ---
% Informações de dados para CAPA e FOLHA DE ROSTO
% ---
\titulo{Auxílio ao resgate de vítimas por robôs em ambientes acidentados}
\autor{Atílio Antônio Dadalto \\ Ícaro Madalena do Nascimento}
\local{Vitória}
\data{2019}
\instituicao{%
  Universidade Federal do Espírito Santo
  \par Departamento de Informática}
\tipotrabalho{Relatório}
\preambulo{Relatório apresentado como requisito parcial para aprovação na disciplina de Programação I, pela Universidade Federal do Espírito Santo.}

\newcommand{\versao}{2.0}
\newcommand{\subtitulo}{Anteprojeto}

% Macros específicas do trabalho.
% Inclua aqui termos que são utilizados muitas vezes e que demandam formatação especial.
% Use sempre \xspace para que o LaTeX inclua espaço em branco após a macro somente quando necessário.
\newcommand{\istar}{\textit{i}$^\star$\xspace}
\newcommand{\java}{Java\texttrademark\xspace}
\newcommand{\latex}{\LaTeX\xspace}

\newfontfamily{\consolas}{Consolas.ttf}

\definecolor{gray}{rgb}{0.5,0.5,0.5}

\lstset{
  language=C,
  basicstyle=\footnotesize\consolas,
%   numbers=left,
%   numberstyle=\tiny,
  frame=tb,
  tabsize=4,
  breaklines=true,
  postbreak=\mbox{\textcolor{red}{$\hookrightarrow$}\space},
  columns=fixed,
  showstringspaces=false,
  showtabs=false,
  keepspaces,
  commentstyle=\color{gray},
  keywordstyle=\color{blue}
}

%%% Configurações finais de aparência. %%%

% Altera o aspecto da cor azul.
\definecolor{blue}{RGB}{41,5,195}

% Informações do PDF.
\makeatletter
\hypersetup{
	pdftitle={\@title}, 
	pdfauthor={\@author},
	pdfsubject={\imprimirpreambulo},
	pdfcreator={LaTeX with abnTeX2},
	pdfkeywords={abnt}{latex}{abntex}{abntex2}{trabalho acadêmico}, 
	colorlinks=true,				% Colore os links (ao invés de usar caixas).
	linkcolor=blue,					% Cor dos links.
	citecolor=blue,					% Cor dos links na bibliografia.
	filecolor=magenta,				% Cor dos links de arquivo.
	urlcolor=blue,					% Cor das URLs.
	bookmarksdepth=4
}
\makeatother

% Espaçamentos entre linhas e parágrafos.
\setlength{\parindent}{1.3cm}
\setlength{\parskip}{0.2cm}


%%% Páginas iniciais do documento: capa, folha de rosto, ficha, resumo, tabelas, etc. %%%

% Inicia o documento.
\begin{document}

% Retira espaço extra obsoleto entre as frases.
\frenchspacing

% Brasão da instituição.
\begin{figure}[h]
  \centering
  \includegraphics[scale=0.055]{figuras/brasao}
\end{figure} 

% Capa do trabalho.
\imprimircapa
\imprimirfolhaderosto

% Índice de capítulos.
\tableofcontents*

%%% Início da parte de conteúdo do documento. %%%

% Marca o início dos elementos textuais.
\clearpage
\textual

% Inclusão dos capítulos.
\chapter*{Introdução}\label{cap-introducao} % Basicamente finalizado

Este trabalho tem por objetivo documentar a estruturação de um projeto de auxílio à tomada de decisões dado um cenário hipotético de um ambiente acidentado. Neste, dados serão extraídos por robôs-célula que investigam o local, que serão processados de forma a aprimorar a atividade de resgate das vítimas que se encontram no ambiente. O plano de testes, presentes no Capítulo~\ref{cap-planejamento-implementacao}, relata a confiabilidade das funções desenvolvidas.

As referências para os métodos de testes, bem como paradigmas aplicativo e recursivo, que serão implementados podem ser encontrados nas notas de aula utilizadas pela disciplina encontradas na página desta em maio de 2019.

Além do plano de testes, no Capítulo~\ref{cap-planejamento-implementacao} também abordamos a forma como planejamos a resolução de cada problema apresentado pelo trabalho, assim como a implementação, de forma geral, das funções mais complexas. Em seguida, o Capítulo~\ref{cap-analise-resultados} realiza uma avaliação da solução final e discorre acerca dos resultados obtidos, também pontuando comentários adicionais.
\chapter{Planejamento e implementação}\label{cap-planejamento-implementacao}

\section{Problema A}\label{problemaA}
\begin{enumerate}[label=\textbf{\alph*)},series=problemas]
\item Calcular a distância percorrida por um determinado robô ao longo do processo de resgate das vítimas. Considere que a distância total percorrida deve ser calculada como a soma de todas as distâncias entre os pontos de passagem do robô.
\end{enumerate}

\noindent \textbf{COMPREENSÃO DO PROBLEMA E PLANEJAMENTO}: deve-se extrair a tupla de um robô da lista de entrada, de acordo com o identificador dado. Em seguida, calcular a distância total que o robô percorreu e retornar esse valor.

Sendo assim, podemos criar as seguintes funções para solucionar o problema:
\begin{enumerate}
    \item Função para somar todos os elementos de uma lista
    \item Função para cálculo de distância no plano cartesiano
    \item Função para extrair uma tupla de uma lista dado um elemento identificador da tupla
    \item Função para calcular a distância entre os pontos percorridos pelo robô, por ordem de instante\label{itemDistInstante}
    \item Função para ordenar lista por instante (talvez não precise disso, precisamos ver o exemplo de entrada)
\end{enumerate}

As duas primeiras são triviais e foram implementadas diversas vezes durante o curso. Na terceira, utilizamos [...] para retornar a tupla que desejávamos. Para o item \ref{itemDistInstante}, [...].

\section{Problema B}\label{problemaB}
\begin{enumerate}[label=\textbf{\alph*)},resume*=problemas]
\item Determine qual dos robôs apresenta o seu último ponto de passagem no terreno de busca que possui a maior distância em relação à origem. Exiba o caminho percorrido pelo robô e o tempo total do percurso;
\end{enumerate}


% # - Extrair, da lista de entrada, o último ponto de passagem de todos os robôs
% # - Determinar qual robô possui ponto mais longe da origem
% # - Imprimir caminho percorrido pelo robô, determinar tempo do percurso

\section{}
\chapter{Análise  e resultados}
\label{sec-metodo}

\begin{apendicesenv}

\chapter{Planos de testes}\label{apendiceA}

Este apêndice serve como repositório para os planos testes das funções mais importantes implementadas no projeto. Algumas funções assumem certos valores de entrada, de forma a afunilar as possíveis entradas, dado que os valores a serem recebidos podem variar enormemente e que o projeto conta com muitas funções; no entanto, nos contemos a assumir entradas de maneira muito razoável, e, além disso, todos os casos básicos foram tratados devidamente dentro das funções.


Para melhor exemplificação e legibilidade dos testes das funções abaixo, usaremos a lista de dados de entrada ilustrada a seguir:

\begin{lstlisting}
listaRobos = [
    ('robo3', 1, (7, 7), 3), ('robo4', 2, (7, 5), 2),
    ('robo3', 3, (5, 4), 3), ('robo3', 4, (8, 1), 4),
    ('robo4', 5, (4, 5), 3), ('robo5', 6, (7, 7), 4)
]
\end{lstlisting}

\section{\texttt{distancia\_total()}}
A função calcula a distância total entre uma lista de pontos somando a distância entre um ponto e o próximo até o fim da lista.

\begin{center}
    \begin{tabular}{|c|c|}
        \cline{1-2}
        lista & resultado esperado \\ \hline
        [] & None \\ \hline  
        lista com um ponto & None \\ \hline  
        [p1, p2, p3] & distância de p1 até p2 + distância de p2 até p3 \\ \hline
    \end{tabular}
\end{center}

Definindo entradas:

\begin{center}
    \begin{tabular}{|c|c|c|}
        \cline{1-3}
        lista & resultado esperado & resultado obtido  \\ \hline
        [] & None & None \\ \hline
        [(4,4)] & None & None \\ \hline
        [(0,0),(3,4),(5,6)] & 7.82842712474619 & 7.82842712474619 \\ \hline

    \end{tabular}
\end{center}

\section{\texttt{distancia\_total\_robo()}}
A função calcula a distância total percorrida por um robô, ou seja, a soma das distâncias de todos os pontos percorridos por ele.

\begin{center}
    \begin{tabular}{|c|c|}
        \cline{1-2}
        lista & resultado esperado \\ \hline
        [] & None \\ \hline  
        [(a,(p1)),(a,(p2))], 'a' & distância de (0,0) até p1 + distância de p1 até p2 \\ \hline
    \end{tabular}
\end{center}


Definindo entradas:

\begin{center}
    \begin{tabular}{|c|c|c|}
        \cline{1-3}
        lista & resultado esperado & resultado obtido  \\ \hline
        [] & None & None \\ \hline
        [('robo1',2,(4,4),5)], 'robo2' & None & None \\ \hline
        listaRobos, 'robo3' & 17.74768689919494 & 17.74768689919494 \\ \hline

    \end{tabular}
\end{center} 


\section{\texttt{indices\_maximos()}}
A função recebe uma lista numérica e retorna uma lista com os índices de todas as ocorrências do valor máximo na lista.

\begin{center}
    \begin{tabular}{|c|c|}
        \cline{1-2}
        lista & resultado esperado \\ \hline
        [] & [] \\ \hline  
        [a, b, c, d] & índices das ocorrências do maior elemento da lista  \\ \hline
    \end{tabular}
\end{center}

Definindo entradas:

\begin{center}
    \begin{tabular}{|c|c|c|}
        \cline{1-3}
        lista & resultado esperado & resultado obtido  \\ \hline
        [] & [] & [] \\ \hline
        [5] & [0] & [0] \\ \hline
        [1,2,3,10,1,10] & [3,5] & [3,5] \\ \hline

    \end{tabular}
\end{center} 


\section{\texttt{imprime\_robos\_mais\_distantes()}}
A função recebe uma lista de robôs e imprime os robôs que encontram-se mais distantes da origem, seus percursos e tempos de percurso
\begin{center}
    \begin{tabular}{|c|c|}
        \cline{1-2}
        lista & resultado esperado \\ \hline
        [] & None \\ \hline  
        [('b',(10,10),1)] & ['b'],[(0,0),(10,10)],[1] \\ \hline
        [('a',(5,4),1),('b',(10,10),1)] & ['b'],[(0,0),(10,10)],[1]  \\ \hline
    \end{tabular}
\end{center}

Definindo entradas:

\begin{center}
    \begin{tabular}{|c|c|c|}
        \cline{1-3}
        lista & resultado esperado & resultado obtido  \\ \hline
        [] & None & None \\ \hline
        [('robo',1,(10,10),3)] & [('robo',((0,0),(10,10)),1] & [('robo',((0,0),(10,10)),1] \\ \hline
        [('robo',1,(10,10),3),('robo2',1,(2,1),3)] & [('robo',((0,0),(10,10)),3] & [('robo',((0,0),(10,10)),3]  \\ \hline

    \end{tabular}
\end{center}




\section{\texttt{caminhos\_robos\_crescente()}}
A função retorna uma lista de tuplas para cada robô, ordenadas em ordem crescente pela distância total percorrida pelos robôs.

\begin{center}
    \begin{tabular}{|c|c|}
        \cline{1-2}
        Entradas & resultado esperado \\ \hline
        [] &  None\\ \hline
        [(a),(b)] & soma do caminho percorrido a, soma do caminho b  \\ \hline
    \end{tabular}
\end{center}


Definindo entradas (saídas arredondadas para melhor visualização):

\begin{center}
    \begin{tabular}{|c|c|}
        \cline{1-2}
        lista & resultado esperado \\ \hline
        [] & None \\ \hline
        [('robo1',2,(4,4),5)] & [('robo1', 5.6, [(0, 0), (4, 4)])] \\ \hline
        [('robo1',2,(4,4),5),('robo3',7,(5,5),2)] & [('robo1', 5.6, [(0, 0), (4, 4)]), ('robo3', 7, [(0, 0), (5, 5)])] \\ \hline
    \end{tabular}
\end{center} 

\begin{center}
    \begin{tabular}{|c|}
        \cline{1-1}
        resultado obtido \\ \hline
        None \\ \hline
        [('robo1', 5.6, [(0, 0), (4, 4)])] \\ \hline
        [('robo1', 5.6, [(0, 0), (4, 4)]), ('robo3', 7, [(0, 0), (5, 5)])] \\ \hline
    \end{tabular}
\end{center}



\section{\texttt{merge\_ordenada\_tupla()}}
Dadas duas listas ordenadas pelo segundo elemento da tupla, junta-as de forma ordenada.

\begin{center}
    \begin{tabular}{|c|c|}
        \cline{1-2}
        lista & resultado esperado \\ \hline
        [] & [] \\ \hline
        [],[(a,b)] & [(a,b)] \\ \hline
        [(c,d)],[] & [(c,d)] \\ \hline
        [(a,b)],[(c,d)]  & [(a,b),(c,d)] \\ \hline
        [(a,b),(c,d)],[(b,c),(d,e)] & [(a,b),(b,c)(c,d),(d,e)] \\ \hline
    \end{tabular}
\end{center}

Definindo entradas:

\begin{center}
    \begin{tabular}{|c|c|c|}
        \cline{1-3}
        lista & resultado esperado & resultado obtido  \\ \hline
        [],[] & [] & []\\ \hline
        [(1,2)],[(3,4)] & [(1,2),(3,4)] & [(1,2),(3,4)] \\ \hline
        [(1,2),(3,4)],[(5,6),(7,8)] & [(1,2),(3,4),(5,6),(7,8)] & [(1,2),(3,4)],[(5,6),(7,8)] \\ \hline
    \end{tabular}
\end{center}


\section{\texttt{merge\_sort\_tupla()}}
Ordena uma lista de tuplas através do método \textit{merge sort}, tendo como referência o segundo elemento das tuplas.

\begin{center}
    \begin{tabular}{|c|c|}
        \cline{1-2}
        lista & resultado esperado \\ \hline
        [] & [] \\ \hline
        [(a)]  & [(a)] \\ \hline
        [(a,d),(c,b),(b,f)] & [(c,b),(a,d),(b,f)]  \\ \hline
    \end{tabular}
\end{center}


Definindo entradas:

\begin{center}
    \begin{tabular}{|c|c|c|}
        \cline{1-3}
        lista & resultado esperado & resultado obtido  \\ \hline
        [] & [] & [] \\ \hline
        [(1,2)] & [(1,2)] & [(1,2)] \\ \hline
        [(1,4),(2,3)] & [(2,3),(1,4)] & [(2,3),(1,4)] \\ \hline

    \end{tabular}
\end{center}


\section{\texttt{ids\_mais\_vitimas()}}
 A função retorna o(s) id(s) do(s) robô(s) que avistaram o maior número de vítimas.
 
 \begin{center}
    \begin{tabular}{|c|c|}
        \cline{1-2}
        lista & resultado esperado \\ \hline
        [] & None \\ \hline
        um robô  & o id do único robô \\ \hline
        mais de um robô, um robô com mais vítimas avistadas & id do robô com mais vítimas  \\ \hline
        mais de um robô, mais de um robô com mais vítimas & ids dos robôs com mais vítimas \\ \hline
    \end{tabular}
\end{center}

Definindo entradas:

\begin{center}
    \begin{tabular}{|c|c|c|}
        \cline{1-3}
        lista & resultado esperado & resultado obtido  \\ \hline
        [] & None & None \\ \hline
        [('robo1',1,(1,1),2)] & ['robo1'] & ['robo1'] \\ \hline
        [('robo1',1,(1,1),2),('robo2',5,(3,2),5)] & ['robo2'] & ['robo2'] \\ \hline
        [('robo1',1,(1,1),2),('robo2',5,(3,2),2)] & ['robo1', 'robo2'] & ['robo1', 'robo2'] \\ \hline
    \end{tabular}
\end{center} 

\end{apendicesenv}


% Finaliza a parte no bookmark do PDF para que se inicie o bookmark na raiz e adiciona espaço de parte no sumário.
\phantompart

% Marca o fim dos elementos textuais.
\postextual

% Fim do documento.
\end{document}
